\documentclass[conference]{IEEEtran}
\IEEEoverridecommandlockouts

\usepackage{amsmath,amssymb,amsfonts}
\usepackage{algorithmic}
\usepackage{graphicx}
\usepackage{textcomp}
\usepackage{xcolor}
\usepackage{booktabs}
\usepackage{cite}
\usepackage{url}
\usepackage{multirow}
\usepackage{microtype}

\begin{document}

\title{Operating-Condition Normalization in C-MAPSS Turbofan Prognostics: A Controlled Empirical Ablation, Label Scarcity, and Evaluation-Protocol Study}

\author{\IEEEauthorblockN{Mr. Neelabh Sharma}
\IEEEauthorblockA{\textit{Department of Mechanical and Aerospace Engineering} \\
\textit{Prognostics and Health Management Laboratory} \\
Email: neelabh.sharma@research.org}
}

\maketitle

\begin{abstract}
Remaining Useful Life (RUL) estimation for aircraft gas turbine engines requires distinguishing mechanical degradation from operating condition shifts. On the multi-regime NASA C-MAPSS benchmark subsets (FD002 and FD004), per-operating-regime sensor normalization is established prior art, yet its empirical contribution is often confounded with architectural innovations or divergent evaluation protocols. This study presents a controlled empirical benchmark evaluating operating-condition normalization, auxiliary regime representations, label scarcity, and evaluation protocols under fixed architectures. Using unsupervised $K$-Means clustering ($k=6$) fit solely on unlabeled training flight settings, per-regime standardization reduces test Root Mean Squared Error (RMSE) over global standardization by $-6.6\%$ on FD002 ($27.07 \pm 0.63$ vs. $28.98 \pm 1.23$) and by $-13.9\%$ on FD004 ($26.79 \pm 0.68$ vs. $31.13 \pm 1.15$) for a 2-layer Gated Recurrent Unit (GRU). Linear models exhibit comparable improvements ($-10.6\%$ on FD002; $-14.5\%$ on FD004), confirming that the benefit is not an idiosyncratic deep learning artifact. Controlled ablation demonstrates that providing regime information as auxiliary features (regime one-hot indicators or standardized flight settings) reduces test RMSE by at most $0.48$ cycles, leaving a $3.86$ to $4.39$ cycle deficit compared to per-regime normalization ($>5$ pooled standard deviations). Extending training to 100 epochs preserves a $3.58$-cycle gap ($30.23 \pm 0.06$ vs. $26.65 \pm 0.43$), ruling out undertraining. Under label scarcity on FD004, a condition-aware GRU trained on $10\%$ of labeled engines ($31.36 \pm 0.85$) achieves statistical equivalence with a full-fleet globally normalized model ($31.13 \pm 1.15$) under a pre-registered paired Two One-Sided Tests (TOST) margin of $\pm 1.5$ cycles ($p = 0.0198$, $95\%$ CI: $[-0.88, +1.35]$). Finally, post-hoc test ground-truth clipping (Protocol B) lowers nominal RMSE by $11$--$12$ cycles solely due to target truncation while preserving the relative benefit of per-regime normalization ($-24.5\%$).
\end{abstract}

\begin{IEEEkeywords}
Remaining useful life, prognostics and health management, C-MAPSS benchmark, sensor normalization, operating regimes, label scarcity, reproducibility.
\end{IEEEkeywords}

\section{Introduction}
Prognostics and Health Management (PHM) for aircraft gas turbine engines aims to prevent catastrophic in-flight failures and optimize fleet maintenance by predicting Remaining Useful Life (RUL) from multivariate run-to-failure telemetry~\cite{saxena2008cmapss}. Turbofan engines operate across wide dynamic flight envelopes characterized by varying altitudes, flight Mach numbers, and throttle resolvers~\cite{frederick2007cmapss}. These operational variations induce telemetry shifts that are often substantially larger in amplitude than the subtle, monotonic sensor drifts caused by internal mechanical degradation (e.g., high-pressure compressor fouling or turbine blade erosion).

On the Commercial Modular Aero-Propulsion System Simulation (C-MAPSS) benchmark~\cite{saxena2008cmapss}, standardizing sensor channels globally across heterogeneous operating regimes conflates operating state changes with physical degradation. Condition-based sensor normalization—clustering operational flight settings and standardizing telemetry within each discovered regime—was proposed in early benchmark studies~\cite{pasa2019}. However, subsequent literature has frequently combined preprocessing modifications with novel neural network architectures (e.g., deep CNNs~\cite{li2018tcn}, domain-adversarial networks~\cite{da2021transfer, ragab2020adversarial}, or spatio-temporal graph attention networks~\cite{stagnn2024}), different sequence windowing parameters, or conflicting test evaluation protocols~\cite{ramasso2014investigating, gupta2024auditing}. As a result, the standalone contribution of operating-condition normalization remains obscured.

Furthermore, acquiring verified run-to-failure degradation trajectories in operational aviation is economically prohibitive, whereas unlabeled operational telemetry across flight envelopes is abundantly available. Evaluating whether unsupervised flight regime normalization can compensate for scarce run-to-failure labels is of immediate practical relevance for fleet deployment.

This paper provides a controlled empirical study addressing four specific questions:
\begin{enumerate}
    \item \textbf{Isolated Preprocessing Contribution:} Across linear (Ridge), ensemble (Random Forest), and recurrent (GRU) models, what is the exact performance delta between global and per-regime standardization under strictly identical architectures?
    \item \textbf{Normalization vs. Auxiliary Regime Inputs:} Does per-regime standardization provide a unique inductive bias, or can a global neural network match its performance when supplied with explicit regime indicators (one-hot vectors) or continuous operating settings?
    \item \textbf{Label Scarcity and Equivalence:} Can unsupervised regime clustering on the entire fleet compensate for scarce failure labels ($10\%$, $25\%$, $50\%$), and does it satisfy statistical equivalence bounds (TOST)?
    \item \textbf{Protocol and Training Confounders:} How much nominal performance variance is driven by evaluation protocol choices (unclipped Protocol A vs. clipped Protocol B) and training epoch limits?
\end{enumerate}

\section{Related Work}

\subsection{Benchmark Evolution on C-MAPSS}
Since the release of the NASA C-MAPSS dataset by Saxena et al.~\cite{saxena2008cmapss}, data-driven prognostics has evolved through feedforward networks~\cite{heimes2008recurrent}, recurrent networks~\cite{zheng2017long, cho2014gru}, temporal convolutional networks~\cite{li2018tcn}, and attention-based spatio-temporal graph networks~\cite{stagnn2024}. Early benchmark surveys by Ramasso and Saxena~\cite{ramasso2014investigating} highlighted that published comparisons frequently suffered from unreported preprocessing steps, inconsistent validation holdouts, and diverging metric definitions.

\subsection{Operating Condition Normalization}
Subsets FD002 and FD004 simulate flights across six discrete operating regimes. Pasa et al.~\cite{pasa2019} introduced condition-invariant sensor standardization via $K$-Means clustering to isolate degradation indicators. Gupta~\cite{gupta2024auditing} demonstrated that many modern deep architectures reported as state-of-the-art were compared against globally normalized baselines, obscuring whether accuracy gains stemmed from model capacity or regime alignment. In this work, we treat per-regime normalization strictly as established prior art and systematically isolate its empirical effect across multiple inductive biases.

\subsection{Label Scarcity and Semi-Supervised Prognostics}
Because run-to-failure events are rare in fielded fleets, semi-supervised architectures~\cite{listou2020semi} and unsupervised domain adaptation~\cite{da2021transfer, ragab2020adversarial} have been explored. Unsupervised regime clustering exploits the fleet's unlabeled operational telemetry without requiring failure annotations, offering an accessible semi-supervised inductive bias.

\subsection{Evaluation Protocol Divergence}
Heimes~\cite{heimes2008recurrent} introduced a piecewise-linear degradation model that caps maximum training RUL targets (typically at $125$ cycles) to avoid penalizing healthy early-life stages. However, published studies diverge fundamentally on whether test ground-truth targets are evaluated in their raw unclipped form (Protocol A) or post-hoc clipped at $125$ cycles (Protocol B)~\cite{ramasso2014investigating, gupta2024auditing}. This distinction produces nominal metric disparities of over $10$ RMSE cycles on identical predictions.

\section{Benchmark Characteristics and Protocols}

\subsection{C-MAPSS Subset Properties}
The benchmark consists of four subsets (Table~\ref{tab:subsets}). Each flight cycle records three operating settings (altitude, Mach number, throttle resolver angle) and 21 sensor channels. Subsets FD001 and FD003 operate at sea level, while FD002 and FD004 cycle across six distinct operating regimes.

\begin{table}[t]
\caption{Characteristics of C-MAPSS Subsets.}
\label{tab:subsets}
\centering
\begin{tabular}{lcccc}
\toprule
Subset & Regimes & Fault Modes & Train Engines & Test Engines \\
\midrule
FD001 & 1 & 1 (HPC) & 100 & 100 \\
FD002 & 6 & 1 (HPC) & 260 & 259 \\
FD003 & 1 & 2 (HPC, Fan) & 100 & 100 \\
FD004 & 6 & 2 (HPC, Fan) & 249 & 248 \\
\bottomrule
\end{tabular}
\end{table}

\subsection{Test Evaluation Protocols and Metrics}
Prognostic evaluation is conducted on the final observed sensor sequence for each test engine. We audit two protocols:
\begin{itemize}
    \item \textbf{Protocol A (Standard Unclipped):} Predictions $\hat{y}_i$ are evaluated directly against the raw true test ground-truth RUL $y_i$.
    \item \textbf{Protocol B (Clipped Ground Truth):} Predictions $\hat{y}_i$ are evaluated against clipped ground truth $\min(y_i, 125)$, matching the training target ceiling.
\end{itemize}

Model accuracy is quantified via Root Mean Squared Error (RMSE) and the asymmetric NASA PHM 2008 Challenge score~\cite{saxena2008cmapss}:
\begin{equation}
S = \sum_{i=1}^N h(d_i), \quad h(d_i) = \begin{cases} e^{-d_i/13} - 1, & d_i < 0 \\ e^{d_i/10} - 1, & d_i \ge 0 \end{cases}
\end{equation}
where $d_i = \hat{y}_i - y_i$. Positive errors (late predictions, predicting longer life than remains) are penalized more severely than negative errors (early predictions).

\section{Methodology}

\subsection{Unsupervised Operating Regime Clustering}
In FD002 and FD004, $K$-Means clustering ($k=6$, $\text{random\_state}=0$, $10$ initializations) is applied to the three operating settings~\cite{lloyd1982least}. To prevent data leakage, cluster centroids are fit strictly on training flight settings across all training engines without target labels. Test engine cycles are mapped to the nearest centroid. 

Silhouette analysis confirms that $k=6$ produces an acute cluster separation peak ($>0.9996$ on both FD002 and FD004), while cluster inertia collapses from $1.34 \times 10^7$ ($k=2$) to $0.47$ ($k=6$).

\subsection{Sensor Normalization Schemes}
Let $x_{i,t,j}$ denote the raw reading of sensor $j$ on engine $i$ at cycle $t$, and let $c(i,t) \in \{1, \dots, k\}$ denote its assigned regime:
\begin{itemize}
    \item \textbf{Global Standardization (Mode G):} $\tilde{x}_{i,t,j} = (x_{i,t,j} - \mu_j) / \sigma_j$, where $\mu_j$ and $\sigma_j$ are computed globally across all training cycles.
    \item \textbf{Per-Regime Standardization (Mode C):} $\tilde{x}_{i,t,j} = (x_{i,t,j} - \mu_{j,c}) / \sigma_{j,c}$, where $\mu_{j,c}$ and $\sigma_{j,c}$ are computed exclusively on training cycles mapped to regime $c$.
\end{itemize}
Sensors with global training variance below $10^{-10}$ exhibit zero fleet dynamics and are dropped. In Mode C, if a sensor's intra-regime standard deviation falls below $10^{-10}$, it is set to $1.0$ to center without distortion. For single-regime subsets (FD001, FD003), $k=1$, making Mode C mathematically identical to Mode G.

\section{Experimental Configuration}
All neural experiments utilize a 2-layer GRU ($h=64$, dropout $0.2$)~\cite{cho2014gru} processing sliding sequences of $30$ cycles with stride $5$. Training minimizes MSE with targets scaled via $y/125.0$. Optimization uses Adam ($lr = 2\times 10^{-3}$, gradient clipping $1.0$), batch size $512$, and $\text{ReduceLROnPlateau}$ (patience 3, factor 0.5) over a maximum of $40$ epochs with early stopping (patience 8).

Validation splits are strictly engine-level: $20\%$ of engines are held out ($50$ validation engines and $199$ training engines on FD004), eliminating cross-window leakage. Test evaluation uses mini-batches of $2048$ sequences.

\begin{table*}[t]
\caption{Controlled Benchmark Results across NASA C-MAPSS Subsets (Protocol A, Unclipped Test Ground Truth). Values reported as mean $\pm$ standard deviation across seeds 0, 1, 2 for stochastic models. Negative percentages denote error reductions.}
\label{tab:main_results}
\centering
\begin{tabular}{llcccccc}
\toprule
Subset & Metric & Ridge (Global)$^\dagger$ & Ridge (Cond)$^\dagger$ & RF (Global) & RF (Cond) & GRU (Global) & GRU (Cond) \\
\midrule
\multirow{2}{*}{\textbf{FD001}} & RMSE & 18.24 & 18.24 & $18.49 \pm 0.13$ & $18.48 \pm 0.13$ & $15.94 \pm 0.29$ & $15.94 \pm 0.29$ \\
 & NASA & 537.3 & 537.3 & $878.4 \pm 54.4$ & $879.2 \pm 59.0$ & $463.8 \pm 84.9$ & $463.8 \pm 84.9$ \\
\midrule
\multirow{2}{*}{\textbf{FD002}} & RMSE & 32.36 & 28.92 ($-10.6\%$) & $30.89 \pm 0.05$ & $29.93 \pm 0.09$ ($-3.1\%$) & $28.98 \pm 1.23$ & $\mathbf{27.07 \pm 0.63}$ ($\mathbf{-6.6\%}$) \\
 & NASA & 14210.8 & 11045.2 & $14547.0 \pm 743.8$ & $13751.6 \pm 330.0$ & $18033.3 \pm 6705.3$ & $\mathbf{8573.4 \pm 1276.2}$ \\
\midrule
\multirow{2}{*}{\textbf{FD003}} & RMSE & 18.77 & 18.77 & $19.61 \pm 0.23$ & $19.60 \pm 0.26$ & $19.33 \pm 2.89$ & $19.33 \pm 2.89$ \\
 & NASA & 651.5 & 651.5 & $1507.4 \pm 56.7$ & $1501.1 \pm 62.2$ & $1247.7 \pm 557.7$ & $1247.7 \pm 557.7$ \\
\midrule
\multirow{2}{*}{\textbf{FD004}} & RMSE & 35.72 & 30.53 ($-14.5\%$) & $32.04 \pm 0.15$ & $30.32 \pm 0.09$ ($-5.4\%$) & $31.13 \pm 1.15$ & $\mathbf{26.79 \pm 0.68}$ ($\mathbf{-13.9\%}$) \\
 & NASA & 12450.4 & 8120.6 & $9008.6 \pm 230.9$ & $7593.8 \pm 139.5$ & $9214.6 \pm 1550.8$ & $\mathbf{4950.5 \pm 403.8}$ \\
\bottomrule
\end{tabular}
\begin{flushleft}
\footnotesize $^\dagger$Ridge regression is deterministic with zero seed variance. On single-regime subsets (FD001, FD003), Mode C is mathematically identical to Mode G ($k=1$).
\end{flushleft}
\end{table*}

\section{Experimental Results and Analysis}

\subsection{Controlled Main Benchmark Performance}
Table~\ref{tab:main_results} presents benchmark performance across all four subsets under Protocol A. On single-regime controls (FD001 and FD003), per-regime standardization defaults to $k=1$, replicating global scaling ($15.94$ on FD001; $19.33$ on FD003).

On multi-regime subsets:
\begin{itemize}
    \item On \textbf{FD002}, GRU test RMSE decreases from $28.98 \pm 1.23$ to $27.07 \pm 0.63$ ($-6.6\%$), and the NASA score decreases from $18033.3$ to $8573.4$ ($-52.5\%$).
    \item On \textbf{FD004}, GRU test RMSE decreases from $31.13 \pm 1.15$ to $26.79 \pm 0.68$ ($-13.9\%$), and the NASA score drops from $9214.6$ to $4950.5$ ($-46.3\%$).
\end{itemize}

Crucially, linear Ridge regression~\cite{hoerl1970ridge} exhibits substantial error reductions under per-regime standardization: $-10.6\%$ on FD002 ($32.36 \to 28.92$) and $-14.5\%$ on FD004 ($35.72 \to 30.53$). Random Forest~\cite{breiman2001random} similarly improves by $-3.1\%$ and $-5.4\%$. On FD004, a condition-aware Random Forest ($30.32 \pm 0.09$) outperforms a globally normalized deep GRU ($31.13 \pm 1.15$), demonstrating that proper physical state alignment can match or exceed deep representation capacity.

\subsection{Double-Bootstrap and Hypothesis Testing}
We evaluate statistical significance across seed and engine dimensions:
\begin{itemize}
    \item \textbf{Seed-level paired $t$-tests ($n=3$):} On FD004, the paired seed difference is statistically significant ($t = -12.663, p = 0.0062$, Cohen's $d = -7.31$~\cite{cohen1988power}). On FD002, with $n=3$, the paired $t$-test yields $t = -2.544, p = 0.126$ ($d = -1.47$), failing significance at $\alpha = 0.05$ due to constrained degrees of freedom ($\text{df}=2$).
    \item \textbf{Engine-level tests ($n=248$ on FD004, $n=259$ on FD002):} Paired $t$-tests on seed-averaged squared errors are highly significant on both subsets ($t = -2.720, p = 0.00697$ on FD002; $t = -5.936, p = 9.82 \times 10^{-9}$ on FD004). Non-parametric Wilcoxon signed-rank tests confirm significance ($p = 0.00374$ on FD002; $p = 3.50 \times 10^{-9}$ on FD004). Holm-adjusted $p$-values~\cite{holm1979sequential} remain significant ($p_{\text{adj}} = 0.00697$ on FD002; $p_{\text{adj}} = 1.96 \times 10^{-8}$ on FD004).
    \item \textbf{Double-Bootstrap CIs ($B=10{,}000$):} Resampling both test engines and training seed models with replacement~\cite{efron1993bootstrap} accounts simultaneously for test fleet sampling error and training stochasticity. The resulting $95\%$ confidence interval for $\Delta\text{RMSE}$ is $[-3.61, -0.23]$ cycles on FD002 and $[-5.80, -2.89]$ cycles on FD004. For the NASA score, the double-bootstrap $95\%$ CI is $[-26883.6, +154.0]$ on FD002 (spanning zero) and $[-6724.8, -1894.1]$ on FD004.
\end{itemize}

\subsection{Ablation: Normalization vs. Auxiliary Regime Inputs (E1)}
A critical open question is whether per-regime normalization provides a structural benefit that cannot be matched by simply providing regime information as extra input features to a globally normalized network. We evaluated five controlled conditions on FD004 under identical training hyperparameters (Table~\ref{tab:e1_summary}):
\begin{itemize}
    \item \textbf{C1:} Global normalization ($24$ features).
    \item \textbf{C2:} Per-regime normalization ($24$ features).
    \item \textbf{C3:} Global normalization $+$ one-hot regime indicators ($30$ features).
    \item \textbf{C4:} Global normalization $+$ standardized raw operating settings ($27$ features).
    \item \textbf{C5:} Per-regime normalization $+$ one-hot regime indicators ($30$ features).
\end{itemize}

\begin{table}[t]
\caption{Ablation of Auxiliary Regime Information on FD004 (Seeds 0, 1, 2). Test metrics under Protocol A.}
\label{tab:e1_summary}
\centering
\begin{tabular}{llccc}
\toprule
Cond. & Description & Val RMSE & Test RMSE & NASA Score \\
\midrule
\textbf{C1} & Global Normalization & $18.26 \pm 1.35$ & $31.13 \pm 1.15$ & $9214.6 \pm 1550.8$ \\
\textbf{C2} & Per-Regime ($k=6$) & $\mathbf{13.94 \pm 0.83}$ & $\mathbf{26.79 \pm 0.68}$ & $\mathbf{4950.5 \pm 403.8}$ \\
\textbf{C3} & Global $+$ One-Hot & $17.72 \pm 1.09$ & $30.65 \pm 0.25$ & $9835.3 \pm 1666.2$ \\
\textbf{C4} & Global $+$ Raw Settings & $18.26 \pm 1.39$ & $31.18 \pm 0.86$ & $9372.0 \pm 980.7$ \\
\textbf{C5} & Per-Regime $+$ One-Hot & $13.87 \pm 0.57$ & $27.02 \pm 0.18$ & $5116.5 \pm 139.1$ \\
\bottomrule
\end{tabular}
\end{table}

The empirical evidence is definitive:
\begin{itemize}
    \item Appending one-hot regime indicators (C3) reduces global test RMSE by only $0.48$ cycles ($31.13 \to 30.65$), remaining $+3.86$ cycles ($7.54$ pooled standard deviations) worse than per-regime normalization (C2).
    \item Supplying continuous standardized flight settings (C4) yields no improvement ($31.18 \pm 0.86$), remaining $+4.39$ cycles ($5.68$ pooled standard deviations) worse than C2.
    \item Combining per-regime normalization with one-hot indicators (C5) yields $27.02 \pm 0.18$, statistically indistinguishable from C2 ($+0.23$ cycles).
\end{itemize}
These results prove that standard recurrent units cannot autonomously synthesize the condition-specific centering and scaling transformations from auxiliary input features.

\subsection{Sensitivity to Training Epoch Cap (Tier 1A)}
Because globally normalized models frequently reached the 40-epoch cap during pilot runs, we audited whether global models were prematurely truncated. We extended the maximum epoch cap to 100 epochs with identical early stopping (patience 8).

\begin{table}[t]
\caption{Extended Training Sensitivity (100-Epoch Cap) on FD004.}
\label{tab:tier1a}
\centering
\begin{tabular}{lcccc}
\toprule
Condition & Seed & Stopped Epoch & Val RMSE & Test RMSE (A) \\
\midrule
\multirow{3}{*}{\textbf{C1} (Global)} & 0 & 79 & 17.60 & 30.25 \\
 & 1 & 54 & 17.11 & 30.27 \\
 & 2 & 100 & 17.52 & 30.16 \\
\cline{2-5}
\multicolumn{2}{l}{\textbf{C1 Mean $\pm$ Std}} & $77.7 \pm 23.0$ & $17.41 \pm 0.26$ & $\mathbf{30.23 \pm 0.06}$ \\
\midrule
\multirow{3}{*}{\textbf{C2} (Per-Regime)} & 0 & 39 & 14.18 & 26.33 \\
 & 1 & 48 & 13.01 & 26.48 \\
 & 2 & 76 & 13.92 & 27.14 \\
\cline{2-5}
\multicolumn{2}{l}{\textbf{C2 Mean $\pm$ Std}} & $54.3 \pm 19.4$ & $13.70 \pm 0.61$ & $\mathbf{26.65 \pm 0.43}$ \\
\bottomrule
\end{tabular}
\end{table}

Table~\ref{tab:tier1a} details the outcome. While extended training allows global models to reach their validation plateau (average stopped epoch $77.7$), reducing test RMSE from $31.13$ to $30.23$ ($-0.90$ cycles), condition-aware models also benefit ($26.79 \to 26.65$). The performance gap between global and per-regime models remains $\mathbf{3.58\text{ cycles}} \ge 3.0$ cycles ($30.23 - 26.65$). This confirms that the advantage of per-regime normalization is robust and not an artifact of undertraining.

\subsection{Label Scarcity Benchmark and Equivalence Testing}
Table~\ref{tab:label_efficiency} presents performance under label constraints on FD004 ($10\%$, $25\%$, $50\%$, $100\%$ labeled training engines). Per-regime normalization consistently reduces test RMSE by $12.6\%$ to $16.3\%$ across all fractions.

\begin{table}[t]
\caption{Label Scarcity Benchmark on Multi-Regime FD004.}
\label{tab:label_efficiency}
\centering
\begin{tabular}{lccc}
\toprule
Fraction (Engines) & Global RMSE & Cond-Aware RMSE & $\Delta$ RMSE (\%) \\
\midrule
10\% ($n=20$) & $37.49 \pm 0.79$ & $\mathbf{31.36 \pm 0.85}$ & $\mathbf{-16.3\%}$ \\
25\% ($n=50$) & $35.38 \pm 0.64$ & $\mathbf{30.90 \pm 0.33}$ & $\mathbf{-12.6\%}$ \\
50\% ($n=100$) & $33.24 \pm 1.03$ & $\mathbf{29.05 \pm 0.80}$ & $\mathbf{-12.6\%}$ \\
100\% ($n=199$) & $31.13 \pm 1.15$ & $\mathbf{26.79 \pm 0.68}$ & $\mathbf{-13.9\%}$ \\
\bottomrule
\end{tabular}
\end{table}

We evaluated formal statistical equivalence using Two One-Sided Tests (TOST)~\cite{schuirmann1987tost, lakens2017tost} under a pre-specified margin $\delta = 1.5$ cycles ($\approx 4.8\%$ of baseline RMSE):
\begin{itemize}
    \item Comparing condition-aware GRU at $10\%$ labels ($31.36 \pm 0.85$) against full-fleet global GRU at $100\%$ labels ($31.13 \pm 1.15$), the mean paired difference across seeds is $+0.23$ cycles ($95\%$ CI: $[-0.88, +1.35]$). Because the entire confidence interval falls strictly within $[-1.5, +1.5]$, the paired TOST test rejects non-equivalence ($p_{\text{TOST}} = 0.0198 < 0.05$). Under an unpaired Welch TOST, $p = 0.1031$ ($\text{df}=3.69$), reflecting power constraints at $n=3$ seeds.
    \item At $25\%$ labels ($30.90 \pm 0.33$), the observed difference is $-0.23$ cycles.
\end{itemize}
These results indicate that unsupervised regime normalization on the full fleet provides an inductive bias that rivals the benefit of acquiring substantial additional failure annotations.

\subsection{Degradation Stage and Error Bin Decomposition}
To evaluate where error reductions occur, test engines were partitioned into true-RUL degradation bins: $0$--$30$, $30$--$60$, $60$--$125$, and $>125$ cycles (Table~\ref{tab:bin_decomp}).

\begin{table}[t]
\caption{Degradation Bin Decomposition on FD004 (Seed 0). Share denotes proportion of total squared error reduction.}
\label{tab:bin_decomp}
\centering
\begin{tabular}{lcccc}
\toprule
True RUL Bin & Global RMSE & Cond RMSE & $\Delta$ RMSE & Share of $\Delta$SE \\
\midrule
$0$--$30$ cycles & 15.02 & 7.35 & $-7.67$ & $14.66\%$ \\
$30$--$60$ cycles & 23.08 & 16.03 & $-7.05$ & $17.21\%$ \\
$60$--$125$ cycles & 20.95 & 18.03 & $-2.92$ & $16.98\%$ \\
$>125$ cycles & 49.52 & 44.58 & $-4.94$ & $51.15\%$ \\
\midrule
\textbf{Total Fleet} & 30.83 & 26.33 & $-4.50$ & $100.0\%$ \\
\bottomrule
\end{tabular}
\end{table}

The total squared error reduction across the 248 test engines is $60{,}899.7$. Crucially, near-failure engines ($0$--$60$ cycles) account for $31.87\%$ of total squared error reduction ($14.66\%$ in $0$--$30$ and $17.21\%$ in $30$--$60$), cutting late-life RMSE nearly in half ($15.02 \to 7.35$). The early-life bin ($>125$ cycles) accounts for $51.15\%$ of squared error reduction, driven by small cycle adjustments across a wide error baseline resulting from the training target ceiling.

\subsection{Evaluation Protocol Divergence (Protocol A vs. B)}
Table~\ref{tab:protocols} summarizes the impact of test ground-truth clipping. Under Protocol B (clipping ground truth at $125$), nominal RMSE drops by $11.8$ cycles on FD004 (from $26.79$ to $14.95$). This metric decrease is purely an artifact of target truncation and does not indicate superior prognostic ability. However, the relative benefit of per-regime normalization is robustly preserved ($-24.5\%$ on FD004; $-13.7\%$ on FD002), confirming that normalization gains hold across both evaluation conventions.

\begin{table}[t]
\caption{Protocol Comparison: Standard (A) vs. Clipped (B) Test Ground Truth.}
\label{tab:protocols}
\centering
\begin{tabular}{llcc}
\toprule
Subset & Method & Protocol A RMSE & Protocol B RMSE \\
\midrule
\multirow{2}{*}{\textbf{FD001}} & GRU Global & $15.94 \pm 0.29$ & $14.66 \pm 0.48$ \\
 & GRU Cond & $15.94 \pm 0.29$ & $14.66 \pm 0.48$ \\
\midrule
\multirow{2}{*}{\textbf{FD002}} & GRU Global & $28.98 \pm 1.23$ & $16.75 \pm 0.16$ \\
 & GRU Cond & $\mathbf{27.07 \pm 0.63}$ & $\mathbf{14.46 \pm 0.19}$ ($-13.7\%$) \\
\midrule
\multirow{2}{*}{\textbf{FD003}} & GRU Global & $19.33 \pm 2.89$ & $18.16 \pm 2.74$ \\
 & GRU Cond & $19.33 \pm 2.89$ & $18.16 \pm 2.74$ \\
\midrule
\multirow{2}{*}{\textbf{FD004}} & GRU Global & $31.13 \pm 1.15$ & $19.81 \pm 1.13$ \\
 & GRU Cond & $\mathbf{26.79 \pm 0.68}$ & $\mathbf{14.95 \pm 0.96}$ ($-24.5\%$) \\
\bottomrule
\end{tabular}
\end{table}

\subsection{Mechanistic Variance Decomposition}
We quantified the proportion of sensor variance explained by operating regimes ($\eta^2 = SS_{\text{between}} / SS_{\text{total}}$). In raw telemetry, flight regimes account for $99.19\%$ of sensor variance on FD002 and $99.25\%$ on FD004 across active sensors. Global standardization leaves between-regime variance completely intact ($\eta^2 = 0.992$). In contrast, per-regime standardization centers each sensor within its regime, driving between-regime variance across all active dynamic sensors to machine precision ($\eta^2 < 10^{-24}$).

\section{Discussion and Threats to Validity}
These findings demonstrate that operating-condition normalization is an essential baseline control in multi-regime prognostics. Studies proposing complex architectures on FD002/FD004 must benchmark against a per-regime normalized baseline to avoid attributing preprocessing effects to architectural capacity.

\textbf{Threats to Validity:}
\begin{itemize}
    \item \textbf{Discrete Flight Setpoints:} C-MAPSS simulates operations at six discrete flight regimes. Operational aircraft experience continuous transitions across flight envelopes, which may require continuous normalization (e.g., kernel regression or spline baselines).
    \item \textbf{Single Benchmark Fleet Split:} NASA C-MAPSS provides a single fixed fleet partition, precluding nested outer cross-validation across independent aircraft operators.
\end{itemize}

\section{Conclusion}
This study presented a controlled empirical ablation of per-regime sensor normalization on the NASA C-MAPSS benchmark. Operating-condition normalization via unsupervised $K$-Means clustering ($k=6$) reduces test RMSE by $-6.6\%$ on FD002 and $-13.9\%$ on FD004 across recurrent networks, with corresponding reductions of $-10.6\%$ and $-14.5\%$ on linear Ridge regression. We showed that auxiliary regime representations (one-hot indicators or continuous operating settings) fail to reproduce this benefit ($>5$ pooled standard deviations worse), and that extended training to 100 epochs preserves a $3.58$-cycle advantage. Under label scarcity, per-regime normalization enables models trained on $10\%$ of labeled engines to match full-fleet global normalization within pre-registered equivalence bounds ($p_{\text{TOST}} = 0.0198$). Finally, protocol clipping reduces nominal RMSE by $11$--$12$ cycles without reflecting algorithmic progress, underscoring the imperative for rigorous, transparent evaluation standards in prognostic research.

\section*{Reproducibility Statement}
All code, clustering routines, raw prediction logs, and execution scripts are organized in the project repository with verified SHA-256 checksums to ensure exact replication.

\bibliographystyle{IEEEtran}
\bibliography{references}

\end{document}
